\documentclass[a4paper, serif, 10pt]{moderncv}

%% moderncv
% color and style
\moderncvcolor{black}
\moderncvstyle{banking}

% renew moderncv command to adapt for CJK colon
\renewcommand*{\cvitem}[3][.25em]{%
  \ifstrempty{#2}{}{\hintstyle{#2}：}{#3}%
  \par\addvspace{#1}}

\renewcommand*{\cvitemwithcomment}[4][.25em]{%
  \savebox{\cvitemwithcommentmainbox}{\ifstrempty{#2}{}{\hintstyle{#2}：}#3}%
  \setlength{\cvitemwithcommentmainlength}{\widthof{\usebox{\cvitemwithcommentmainbox}}}%
  \setlength{\cvitemwithcommentcommentlength}{\maincolumnwidth-\separatorcolumnwidth-\cvitemwithcommentmainlength}%
  \begin{minipage}[t]{\cvitemwithcommentmainlength}\usebox{\cvitemwithcommentmainbox}\end{minipage}%
  \hfill% fill of \separatorcolumnwidth
  \begin{minipage}[t]{\cvitemwithcommentcommentlength}\raggedleft\small\itshape#4\end{minipage}%
  \par\addvspace{#1}}

%% page layout/margins
\usepackage[top=2.5cm, bottom=2.5cm, left=1.5cm, right=1.5cm]{geometry}

\nopagenumbers{}

%% Babel config for English language
\usepackage[english]{babel}

%% fontspec
\usepackage{fontspec}

\IfFontExistsTF{Linux Libertine}{
  \setmainfont[Ligatures={TeX, Common}, Numbers=Lining, ItalicFont=Linux Libertine]{Linux Libertine}
}{}
\IfFontExistsTF{Linux Libertine O}{
  \setmainfont[Ligatures={TeX, Common}, Numbers=Lining, ItalicFont=Linux Libertine O]{Linux Libertine O}
}{}

%% CTeX
% CJK support, used to show CJK characters in the resume
%
% - heading=false: disable ctex heading
% - punct=kaiming: use kaiming punctuations styles for CJK
% - scheme=plain: use plain scheme, do not override \normalsize font size
% - space=auto: space settings for CJK characters
%
% CJK fonts are resolved per locale language, e.g. Noto Serif/Sans CJK JP
% for Japanese resumes, with fallback to the next region if not installed
%
% ref:
% - http://ctan.mirrorcatalogs.com/language/chinese/ctex/ctex.pdf
\usepackage[UTF8, heading=false, punct=kaiming, scheme=plain, space=auto]{ctex}

\IfFontExistsTF{Noto Serif CJK HK}{
  \setCJKmainfont{Noto Serif CJK HK}
}{
\IfFontExistsTF{Noto Serif CJK TC}{
  \setCJKmainfont{Noto Serif CJK TC}
}{}}
\IfFontExistsTF{Noto Sans CJK HK}{
  \setCJKsansfont{Noto Sans CJK HK}
}{
\IfFontExistsTF{Noto Sans CJK TC}{
  \setCJKsansfont{Noto Sans CJK TC}
}{}}

%% line spacing
\usepackage{setspace}
\setstretch{1.125}

%% URL styling - use normal text instead of monospace
\urlstyle{same}

%% Patch \httplink and \httpslink to handle full URLs
% The original moderncv macros always prepend http:// or https:// to URLs.
% This causes issues when the URL already has a protocol (e.g., https://example.com)
% resulting in malformed URLs like https://https://example.com.
% This patch checks if the URL already contains "://" and uses it directly if so.
% Note: We use \str_set:Nx (x-expansion) to expand macros like \@homepage first.
\ExplSyntaxOn
\str_new:N \g_yamlresume_url_str

\RenewDocumentCommand{\httpslink}{O{}m}{
  \str_set:Nx \g_yamlresume_url_str {#2}
  \str_if_in:NnTF \g_yamlresume_url_str {://}
    { \url{#2} }
    { \url{https://#2} }
}

\RenewDocumentCommand{\httplink}{O{}m}{
  \str_set:Nx \g_yamlresume_url_str {#2}
  \str_if_in:NnTF \g_yamlresume_url_str {://}
    { \url{#2} }
    { \url{http://#2} }
}
\ExplSyntaxOff

\name{陳家豪}{}
\title{資深數據工程師}
\phone[mobile]{+852 9123 4567}
\email{kaho.chan@example.com}
\homepage{https://kahochan.dev}

\address{中國香港，香港島，香港，中環皇后大道中99號}{}{}

\extrainfo{{\small \faGithub}\ \href{https://github.com/kahochan-data}{@kahochan-data} {} {} {} • {} {} {} 
{\small \faLinkedin}\ \href{https://www.linkedin.com/in/kahochan-data}{@kahochan-data}}

\begin{document}

\maketitle

\section{簡介}

\cvline{}{資深數據工程師，擁有超過八年設計與建置現代數據平台嘅經驗，專注於即時數據管道、數據倉儲同雲端數據架構。擅長將複雜業務需求轉化為可擴展、可靠且易維護嘅數據解決方案。
%
\begin{itemize}
\item 精通 Python、SQL、Spark 及 dbt，熟悉 Kafka、Airflow 同 Snowflake 等現代數據工具鏈
\item 曾領導跨職能團隊完成多個大型數據遷移及實時分析項目，平均延遲降低 40\%
\item 重視數據品質、可觀測性同成本效益，致力推動數據治理最佳實踐
\end{itemize}}
\section{教育背景}

\cventry{2012年9月–2016年6月}
        {學士，計算機科學，成績：3.7}
        {香港科技大學}
        {\url{https://www.ust.hk}}
        {}
        {\begin{itemize}
\item 主修計算機科學，副修統計學，畢業成績優異
\item 畢業論文探討分散式資料處理框架於金融風險分析之應用
\item 曾獲院長嘉許名單及多項獎學金
\end{itemize}
\textbf{課程}：資料結構與演算法、資料庫系統、分散式系統、大數據分析、雲端運算、機器學習導論}
\section{工作經歷}

\cventry{2021年3月至今}
        {高級數據工程師}
        {智慧數據科技有限公司}
        {\url{https://www.smartdata-tech.hk}}
        {}
        {\begin{itemize}
\item 帶領四人數據工程團隊，設計及維護每日處理超過 5 TB 數據嘅雲端數據平台
\item 主導由傳統數據倉儲遷移至 Snowflake 嘅項目，查詢效能提升 3 倍，年度成本節省約 25\%
\item 建立基於 Airflow 同 dbt 嘅 ELT 管道，實現數據血緣追蹤同自動化測試，數據事故減少 60\%
\item 與數據科學團隊緊密合作，交付實時推薦系統所需嘅特徵管道，模型上線時間由數週縮短至數日
\end{itemize}
\textbf{關鍵字}：數據架構、Snowflake、Airflow、dbt、團隊領導}

\cventry{2017年7月–2021年2月}
        {數據工程師}
        {金融雲端服務有限公司}
        {\url{https://www.fincloud.hk}}
        {}
        {\begin{itemize}
\item 負責建置及優化銀行交易數據嘅批處理管道，每日處理超過 2 億筆記錄
\item 使用 Spark 及 Kafka 開發近實時風險監控數據流，將欺詐偵測延遲由小時級降至分鐘級
\item 設計分區及壓縮策略，使數據湖儲存成本降低 35\%，同時維持查詢效能
\item 參與數據治理委員會，制定數據品質指標及元數據管理標準
\end{itemize}
\textbf{關鍵字}：Spark、Kafka、數據湖、風險監控、數據治理}

\cventry{2016年8月–2017年6月}
        {數據分析師}
        {初創分析有限公司}
        {\url{https://www.startup-analytics.hk}}
        {}
        {\begin{itemize}
\item 為零售客戶建立銷售及庫存儀表板，支援每日營運決策
\item 以 Python 及 SQL 自動化報表流程，每週節省約 15 小時人手工作
\item 協助設計 A/B 測試框架，分析促銷活動成效並提出優化建議
\end{itemize}
\textbf{關鍵字}：Python、SQL、數據視覺化、A/B 測試}
\section{語言}

\cvline{粵語}{母語或雙語水平}
\cvline{英語}{完全專業水平 \hfill \textbf{關鍵字}：IELTS 8.0}
\cvline{普通話}{完全專業水平}
\section{專業技能}

\cvline{數據工程}{專家 \hfill \textbf{關鍵字}：Python、SQL、Spark、Kafka、Airflow、dbt}
\cvline{雲端平台}{高級 \hfill \textbf{關鍵字}：AWS、GCP、Snowflake、BigQuery、Docker、Kubernetes}
\cvline{數據治理}{高級 \hfill \textbf{關鍵字}：數據品質、元數據管理、數據血緣、隱私合規、成本優化}
\section{榮譽}

\cventry{2022年12月}
        {智慧數據科技有限公司}
        {年度最佳數據項目}
        {}
        {}
        {表揚在 Snowflake 遷移項目中展現卓越技術領導力，顯著提升平台效能並降低營運成本。}

\cventry{2015年6月}
        {香港科技大學}
        {院長嘉許名單}
        {}
        {}
        {因學術成績優異而獲頒，每學年僅頒授予成績排名前 10\% 的學生。}
\section{證書}

\cventry{2022年8月}
        {Amazon Web Services}
        {AWS Certified Data Analytics – Specialty}
        {\url{https://aws.amazon.com/certification/}}
        {}
        {}

\cventry{2020年11月}
        {Google Cloud}
        {Google Professional Data Engineer}
        {\url{https://cloud.google.com/certification/data-engineer}}
        {}
        {}
\section{出版刊物}

\cventry{2016年6月}
        {分散式資料處理框架於金融風險分析之應用}
        {香港科技大學計算機科學系}
        {\url{https://hdl.handle.net/123456789/1234}}
        {}
        {\begin{itemize}
\item 探討 Apache Spark 及 Flink 於即時風險計算嘅效能比較
\item 提出基於微批次嘅優化策略，將延遲降低約 30\%
\end{itemize}}
\section{項目}

\cventry{2023年1月–2023年8月}
        {基於 Kafka、Spark Structured Streaming 及 Redis 嘅實時欺詐偵測系統}
        {實時欺詐偵測管道}
        {\url{https://github.com/kahochan-data/fraud-pipeline}}
        {}
        {\begin{itemize}
\item 設計及實作每秒處理 5,000 筆交易嘅串流管道，端到端延遲低於 2 秒
\item 整合機器學習模型評分服務，實現自動化風險決策
\item 建立監控及告警機制，確保管道可用性達 99.9\%
\end{itemize}
\textbf{關鍵字}：Kafka、Spark、Redis、實時數據}

\cventry{2022年3月–2022年12月}
        {以 dbt 及 Great Expectations 為核心嘅數據品質測試框架}
        {數據品質自動化框架}
        {\url{https://github.com/kahochan-data/dq-framework}}
        {}
        {\begin{itemize}
\item 提供宣告式數據品質規則，支援欄位級及表級驗證
\item 自動產生數據品質報告並整合至 CI/CD 流程
\item 於三個業務團隊推行，數據事故平均修復時間縮短 50\%
\end{itemize}
\textbf{關鍵字}：dbt、Great Expectations、CI/CD、數據品質}
\section{興趣愛好}

\cvline{運動}{龍舟、遠足、游泳}
\cvline{音樂}{結他、爵士樂}
\section{志願服務}

\cventry{2019年9月至今}
        {導師及講者}
        {香港數據工程師社群}
        {\url{https://www.hkde.community}}
        {}
        {\begin{itemize}
\item 定期舉辦數據工程工作坊，分享雲端數據架構及管道設計經驗
\item 指導初入行工程師，協助他們建立實戰項目及職涯規劃
\item 籌辦年度數據馬拉松，吸引超過 200 名參加者
\end{itemize}}
    
\end{document}